\section*{第 \ref{ch:b3:rings} 章——环与算术}

\begin{solution}{exo:b3:rings:1}
(a) 求值映射 $f \colon \Z[X] \to \Z[\iu]$，$P \mapsto P(\iu)$，
是满射环同态（$a + bX \mapsto a + b\iu$）。核：
用 $P$ 除以\emph{首一}多项式 $X^2 + 1$（在 $\Z[X]$ 中）：$P = (X^2 +
1)Q + (bX + a)$，其中 $a, b \in \Z$；于是 $P(\iu) = a + b\iu = 0$
当且仅当 $a = b = 0$。因此 $\ker f = (X^2+1)$，
由\cref{thm:b3:rings:firstiso}即得结论。

(b) 以 $\R$ 为系数做同样的计算：$\R[X]/(X^2+1)
\cong \C$——这是 $\C$ 的最简洁\emph{构造}。

(c) $X^2 + X + 1$ 在 $\mathbb F_2$ 中无根（$0, 1 \mapsto 1$），
故其次数为 $2$，从而不可约：\omterm{def:b3:rings:ideal}{商环} $\mathbb
F_4 = \mathbb F_2[X]/(X^2+X+1)$ 是域
（见\cref{prop:b3:rings:primemaximal}；$(P)$ 在 $K[X]$ 中极大，
$P$ 不可约时，因 $K[X]$ 为\omterm{def:b3:rings:pidufd}{主理想整环}，\omterm{def:b3:rings:ideal}{理想} $(D) \supseteq
(P)$ 意味着 $D \mid P$）。它的四个元素为 $0, 1, \omega,
\omega + 1$，其中 $\omega = \bar X$，且 $\omega^2 = \omega + 1$。
非零元的乘法表：
\[
\omega \cdot \omega = \omega + 1, \qquad
\omega(\omega + 1) = \omega^2 + \omega = 1, \qquad
(\omega+1)^2 = \omega^2 + 1 = \omega .
\]
非零元构成 $3$ 阶循环群，由
$\omega$ 生成。
\end{solution}

\begin{solution}{exo:b3:rings:2}
(a) 设 $p$ 为\omterm{def:b3:rings:pidufd}{唯一分解整环}中的\omterm{def:b3:rings:divisibility}{不可约元}，且 $p \mid ab$，记 $ab =
pc$，其中 $a, b \neq 0$（否则结论平凡）。若 $a$ 或 $b$ 为单位，
则 $p$ 整除另一者。否则将 $a$、$b$ 与 $c$ 分解为
\omterm{def:b3:rings:divisibility}{不可约元}之积，于是 $ab$ 的两种分解：
\[
(\text{ 的因子 } a)(\text{ 的因子 } b) = p \cdot
(\text{ 的因子 } c),
\]
在不计次序与\omterm{def:b3:rings:divisibility}{相伴}的意义下必定一致：$p$ 与某个
$a$ 或 $b$ 的不可约因子\omterm{def:b3:rings:divisibility}{相伴}，从而整除它。

(b) 设 $A$ 为有限整环，$x \neq 0$。映射 $y \mapsto
xy$ 为单射（$xy = xy' \Rightarrow x(y - y') = 0 \Rightarrow
y = y'$），从而（因 $A$ 有限）为满射，故存在满足 $1 = xy$ 的 $y$。

(c) 若 $\mathfrak p$ 是有限环 $A$ 中的\omterm{def:b3:rings:primemaximal}{素理想}，则
$A/\mathfrak p$ 为有限整环，由 (b) 知它是域，故
$\mathfrak p$ 为\omterm{def:b3:rings:primemaximal}{极大理想}（见\cref{prop:b3:rings:primemaximal}）。
\end{solution}

\begin{solution}{exo:b3:rings:3}
范数：$N(3) = 9$，$N(1 \pm \iu\sqrt5) = 6$，$N(2 \pm \iu\sqrt5) =
9$。方程 $x^2 + 5y^2 = 2$ 与 $x^2 + 5y^2 = 3$ 均无
整数解，故无元素的范数为 $2$ 或 $3$。
$3$ 的真因式分解需要两个范数为 $3$ 的因子：
不可能——$3$ 不可约。而 $1
\pm \iu\sqrt5$（范数为 $6$）的真因式分解需要范数分别为 $2, 3$ 的因子：
不可能。对 $2 \pm \iu\sqrt5$ 同理（范数 $9$，
因子须范数为 $3$）。于是
\[
9 = 3 \cdot 3 = (2 + \iu\sqrt5)(2 - \iu\sqrt5),
\]
两种不可约分解。它们确实
不同：单位为 $\pm 1$（范数 $1$），且 $2 \pm \iu\sqrt5
\neq \pm 3$。因此唯一性不成立，而分解的存在性
在 $\Z[\iu\sqrt5]$ 中成立
（见\cref{exo:b3:rings:10}(c)）：此处的非唯一分解纯粹是唯一性失效。
（与此相符，见\cref{prop:b3:rings:primeirred}：
这些\omterm{def:b3:rings:divisibility}{不可约元}都不是\omterm{def:b3:rings:divisibility}{素元}。）
\end{solution}

\begin{solution}{exo:b3:rings:4}
(a) 设 $a, b \in \Z[\iu]$，$b \neq 0$，且 $a/b = x + \iu y \in
\C$。取整数 $m, n$，使得 $\abs{x - m} \leq \frac12$，
$\abs{y - n} \leq \frac12$，并令 $q = m + \iu n$，$r = a - bq$。
于是
\[
N(r) = N(b)\,\abs*{\tfrac ab - q}^2
\leq N(b)\Bigl(\tfrac14 + \tfrac14\Bigr) = \tfrac{N(b)}2 < N(b).
\]
因此 $N$ 为欧几里得函数（$N(r) < N(b)$ 或 $r = 0$）。

(b) 若 $uv = 1$，则 $N(u)N(v) = 1$，由 $N(u) \in \N$ 知 $N(u) =
1$，即 $x^2 + y^2 = 1$：$u \in \{\pm 1, \pm\iu\}$；反之
这些均为单位。

(c) 对 $\Z[\iu\sqrt2]$，同样的取整给出 $\abs{a/b - q}^2
\leq \frac14 + \frac{2}4 = \frac34 < 1$，故为\omterm{def:b3:rings:pidufd}{欧几里得整环}；单位由 $x^2 +
2y^2 = 1$ 得 $\pm 1$。对 $\Z[\iu\sqrt5]$，上界变为
$\frac14 + \frac54 = \frac32 > 1$，取整论证失效——
且必然失效，因 $\Z[\iu\sqrt5]$ 甚至不是\omterm{def:b3:rings:pidufd}{唯一分解整环}
（见\cref{exo:b3:rings:3}），而\omterm{def:b3:rings:pidufd}{欧几里得整环}必为\omterm{def:b3:rings:pidufd}{唯一分解整环}
（见\cref{thm:b3:rings:euclideanpid,thm:b3:rings:pidufd}）。
\end{solution}

\begin{solution}{exo:b3:rings:5}
(a) 若 $x^n = 0$：
\[
(1 + x)\bigl(1 - x + x^2 - \dots + (-1)^{n-1}x^{n-1}\bigr)
= 1 + (-1)^{n-1}x^n = 1 .
\]
(b) 在整环中，$\deg(PQ) = \deg P + \deg Q$；由 $PQ = 1$ 推出
$\deg P = \deg Q = 0$，且 $P, Q \in A^\times$：$A[X]^\times =
A^\times$。在 $\Z/4\Z$ 上，$(1 + 2X)^2 = 1 + 4X + 4X^2 = 1$，故 $1
+ 2X$ 是次数为 $1$ 的单位（此处 $2$ 幂零；比较
(a)）。

(c) 由 $e^2 = e$ 得 $e(e - 1) = 0$，故在整环
中 $e \in \{0, 1\}$。若 $x^n = 0$，$n \geq 1$ 取最小，且 $x \ne 0$，则
$n \geq 2$，且 $x \cdot x^{n-1} = 0$，两个因子均非零：
矛盾。
\end{solution}

\begin{solution}{exo:b3:rings:6}
(a) $K[X,Y]/(X,Y) \cong K$（在 $(0,0)$ 处求值）：为域，故
$(X,Y)$ 为\omterm{def:b3:rings:primemaximal}{极大理想}。若 $(X, Y) = (P)$，由 $P \mid X$ 可得（比较
$Y$ 的次数）$P \in K[X]$，再由 $P \mid Y$ 得 $P \in K$；$P =
0$ 不可能，$P \in K^\times$ 则给出 $(P) = K[X,Y]$，
与真\omterm{def:b3:rings:ideal}{理想}相矛盾（$K[X,Y]/(X,Y) \cong K \neq 0$）。故
$(X,Y)$ 不是主\omterm{def:b3:rings:ideal}{理想}。

(b) 求值映射 $P(X, Y) \mapsto P(X, X^2)$ 将 $K[X,Y]$ 满射到
$K[X]$；其核为 $(Y - X^2)$，以关于 $Y$ 首一的
多项式 $Y - X^2$ 作带余除法，$P = (Y - X^2)Q + R(X)$，于是 $P(X, X^2) =
R(X)$。故 $K[X,Y]/(Y - X^2) \cong K[X]$——即抛物线的
坐标环，与直线的坐标环同构。

求值映射 $P(X, Y) \mapsto P(X, X^{-1})$ 将 $K[X, Y]$ 满射到
洛朗多项式环 $K[X, X^{-1}]$。其核包含
$(XY - 1)$；反之，模 $XY - 1$ 时，每个剩余类都有一个
代表元 $R = \sum_{n \geq 0} a_nX^n + \sum_{m \geq 1} b_m
Y^m$（反复将每个乘积 $XY$ 换为 $1$），且 $R(X,
X^{-1}) = \sum a_n X^n + \sum b_m X^{-m} = 0$ 蕴含所有 $a_n =
b_m = 0$。于是 $K[X, Y]/(XY - 1) \cong K[X, X^{-1}]$——即
双曲线的坐标环：去掉一点的直线。

(c) $(Y - X^2)$ 为\omterm{def:b3:rings:primemaximal}{素理想}（\omterm{def:b3:rings:ideal}{商环} $K[X]$ 为整环），但不是
\omterm{def:b3:rings:primemaximal}{极大理想}（$K[X]$ 不是域；具体地，$(Y - X^2) \subsetneq
(Y - X^2,\, X) \subsetneq K[X,Y]$）。
\end{solution}

\begin{solution}{exo:b3:rings:7}
\emph{$X^5 - 12X^3 + 36X - 12$}：取 $p = 3$，由\omterm{thm:b3:rings:criteria}{艾森斯坦判别法}（$3 \mid
12, 36, 12$；$9 \nmid 12$；$3 \nmid 1$），故不可约。（当 $p = 2$
时\omterm{thm:b3:rings:criteria}{艾森斯坦判别法}失效：$4 \mid 12$。）

\emph{$X^4 + X + 1$}：模 $2$ 约化，$\mathbb F_2$ 中无根；
$\mathbb F_2$ 上唯一的不可约二次多项式是 $X^2 + X +
1$，而 $(X^2+X+1)^2 = X^4 + X^2 + 1 \neq X^4 + X + 1$。故 $X^4 +
X + 1$ 在 $\mathbb F_2$ 上不可约，从而在 $\Q$ 上
（见\cref{thm:b3:rings:criteria}(1)；它首一）。

\emph{$X^4 + 4$}：可约——即索菲·热尔曼恒等式，
$X^4 + 4 = (X^2 - 2X + 2)(X^2 + 2X + 2)$.

\emph{$X^4 + 1$}：作平移，$(X+1)^4 + 1 = X^4 + 4X^3 + 6X^2 + 4X +
2$：对 $2$ 用\omterm{thm:b3:rings:criteria}{艾森斯坦判别法}。$X^4+1$ 的分解经平移后
得 $(X+1)^4 + 1$ 的一个分解，故不可约。

\emph{$X^3 - X - 1$}：三次多项式在 $\Q$ 上可约当且仅当它有一个
有理根；首一整系数多项式的有理根是
能整除常数项的整数（有理根定理：若
$(p/q)$ 为既约分数且是根，则 $q \mid 1$，$p \mid -1$），而
$\pm 1$ 都不是根（代入值分别为 $-1$ 与 $-1$），故不可约。
\end{solution}

\begin{solution}{exo:b3:rings:8}
(a) 设 $\gcd(m, n) = 1$，由\cref{thm:b3:rings:crt}可得环同构
$\Z/mn\Z \cong \Z/m\Z \times \Z/n\Z$。乘积环中的元素为单位，
当且仅当两个分量皆为单位，从而
$(\Z/mn\Z)^\times \cong (\Z/m\Z)^\times \times (\Z/n\Z)^\times$
$\varphi(mn) = \varphi(m)\varphi(n)$。对素数幂，
$\Z/p^k\Z$ 中的非单位即 $p$ 的倍数所在的剩余类：
$\varphi(p^k) = p^k - p^{k-1}$。由此
\[
\varphi(n) = \prod_i \bigl(p_i^{\alpha_i} -
p_i^{\alpha_i-1}\bigr) = n \prod_{p \mid n}\Bigl(1 -
\frac1p\Bigr).
\]

(b) 记 $M = 7 \cdot 11 \cdot 13 = 1001$。幂等元：$e_1 \equiv
(1, 0, 0)$：$143 = 11\cdot13 \equiv 3 \pmod 7$，且 $3 \cdot 5
\equiv 1$，故 $e_1 = 143 \cdot 5 = 715$。$e_2$：$91 \equiv 3
\pmod{11}$，$3 \cdot 4 \equiv 1$，故 $e_2 = 91\cdot4 = 364$。$e_3$：
$77 \equiv -1 \pmod{13}$，故 $e_3 = 77 \cdot 12 = 924$。于是
\[
x \equiv 2\,e_1 + 5\,e_2 + 1\,e_3 = 1430 + 1820 + 924 = 4174
\equiv 170 \pmod{1001},
\]
事实上 $170 = 24\cdot7 + 2 = 15\cdot11 + 5 = 13\cdot13 + 1$。
\end{solution}

\begin{solution}{exo:b3:rings:9}
(a) 若 $x^n = 0$ 且 $y^m = 0$，则 $(x+y)^{n+m}$ 的二项展开式
中每一项 $x^iy^j$ 满足 $i + j = n + m$，于是 $i
\geq n$ 或 $j \geq m$：每一项皆为零，故 $x + y$ 为
幂零元；$(ax)^n = a^nx^n = 0$，故 $\operatorname{Nil}(A)$ 为一个
\omterm{def:b3:rings:ideal}{理想}。若 $\mathfrak p$ 为\omterm{def:b3:rings:primemaximal}{素理想}，且 $x^n = 0 \in \mathfrak p$，
对 $n$ 归纳即得 $x \in \mathfrak p$（$x \cdot x^{n-1} \in
\mathfrak p$）。

(b) 设 $a \notin \operatorname{Nil}(A)$，记 $S = \{a^n : n \geq
1\}$，故 $0 \notin S$。记 $\mathcal E$ 为与 $S$ 不交的\omterm{def:b3:rings:ideal}{理想}所成的集合，
它含 $(0)$，且是归纳的（
与 $S$ 不交的\omterm{def:b3:rings:ideal}{理想}链之并仍为与 $S$ 不交的\omterm{def:b3:rings:ideal}{理想}），由佐恩引理
得 $\mathfrak p \in \mathcal E$ 为极大元。$\mathfrak p$ 为
真\omterm{def:b3:rings:ideal}{理想}（$a \notin \mathfrak p$，因 $a \in S$）。验证素性：设
$x, y \notin \mathfrak p$。由极大性，$\mathfrak p + (x)$ 与
$\mathfrak p + (y)$ 均与 $S$ 相交：$a^m \in \mathfrak p + (x)$，$a^n
\in \mathfrak p + (y)$。相乘得 $a^{m+n} \in \mathfrak p +
(xy)$。若 $xy \in \mathfrak p$，则 $a^{m+n} \in \mathfrak p \cap
S$，矛盾。故 $xy \notin \mathfrak p$——即素性的
逆否形式。因此每个非幂零元都不属于某个\omterm{def:b3:rings:primemaximal}{素
理想}；结合 (a)，
$\operatorname{Nil}(A) = \bigcap_{\mathfrak p} \mathfrak p$.
\end{solution}

\begin{solution}{exo:b3:rings:10}
(a) 升链 $\ker f \subseteq \ker f^2 \subseteq \cdots$
稳定（见\cref{prop:b3:rings:noethacc}）：$\ker f^n = \ker
f^{n+1}$，对某个 $n$ 成立。设 $x \in \ker f$。由 $f$（从而 $f^n$）
为满射，$x = f^n(y)$，对某个 $y$；于是 $f^{n+1}(y) = f(x) =
0$，故 $y \in \ker f^{n+1} = \ker f^n$，即 $x = f^n(y) = 0$。

(b) $I_n = \{f \in \mathcal C(\intcc01, \R) : f\restriction_{
\intcc0{1/n}} = 0\}$ 是一个\omterm{def:b3:rings:ideal}{理想}，且 $I_n \subseteq I_{n+1}$。
包含关系严格成立：$x \mapsto \max\bigl(0, x -
\frac1{n+1}\bigr)$ 在 $\intcc0{\frac1{n+1}}$ 上为零，而在
$\intcc0{\frac1n}$ 上不为零。无穷严格递增链
与\cref{prop:b3:rings:noethacc}矛盾。

(c) 假设不能分解为\omterm{def:b3:rings:divisibility}{不可约元}之积的非零非单位元所成的集合非空。
对应的\omterm{def:b3:rings:ideal}{理想}族
$\{(a)\}$ 有极大元 $(a)$
（\cref{prop:b3:rings:noethacc}）。元素 $a$ 不是
\omterm{def:b3:rings:divisibility}{不可约元}（\omterm{def:b3:rings:divisibility}{不可约元}本身即是其自身的一个分解），因此 $a = bc$
其中 $b, c$ 为非单位元，$(a) \subseteq (b)$ 是严格的（若 $(a) =
(b)$ 将得 $b = ad$，$a = adc$，于是 $dc = 1$，$c$ 为单位元），
同理 $(a) \subsetneq (c)$。由极大性，$b$ 与 $c$ 都
可分解为\omterm{def:b3:rings:divisibility}{不可约元}之积；将这些因子拼接即得 $a$ 的分解，
与假设矛盾。对 $\Z[\iu\sqrt5]$ 应用此论证——它作为
$\Z[X]$ 的商是诺特的（\cref{thm:b3:rings:hilbert}，$\Z[\iu\sqrt5]
\cong \Z[X]/(X^2+5)$）——这表明其中分解存在；
\cref{exo:b3:rings:3} 已表明失败的是唯一性。
\end{solution}

\begin{solution}{exo:b3:rings:11}
(a) \cref{exo:b3:rings:7}：作平移，并在 $2$ 处用\omterm{thm:b3:rings:criteria}{艾森斯坦判别法}。

(b) 模 $2$：$X^4 + 1 = (X + 1)^4$。现设 $p$ 为奇数。
平方元构成指数为 $2$ 的子群，该子群位于 $(\Z/p\Z)^\times$ 中：
同态 $x \mapsto x^2$ 的核为 $\{\pm 1\}$（恰有两个元素：$X^2
- 1$ 在域中至多有 $2$ 个根，且 $1 \neq -1$，
$p$ 为奇数），所以其像有 $\frac{p-1}2$ 个元素。因此
两个非平方元之积是平方元（在阶为 $2$ 的\omterm{thm:b3:groups:quotient}{商群}中，
$\overline{xy} = \bar x\bar y$）。故\emph{$-1$、$2$、$-2$ 中
至少有一个是模 $p$ 的平方}（若 $-1$ 与 $2$ 均非平方，
则 $-2 = (-1)\cdot 2$ 为平方）。各情形下 $X^4 + 1$ 模 $p$ 可分解：
\begin{itemize}
  \item $-1 = c^2$: \quad $X^4 + 1 = X^4 - c^2 = (X^2 - c)(X^2 +
        c)$;
  \item $2 = c^2$: \quad $(X^2 + cX + 1)(X^2 - cX + 1) = X^4 + (2
        - c^2)X^2 + 1 = X^4 + 1$;
  \item $-2 = c^2$: \quad $(X^2 + cX - 1)(X^2 - cX - 1) = X^4 -
        (c^2 + 2)X^2 + 1 = X^4 + 1$.
\end{itemize}
故 $X^4+1$ 模每个素数都可约，但在
$\Q$ 上不可约：约化判别法（\cref{thm:b3:rings:criteria}(1)）
能检测不可约性，但其失效并不说明任何问题。

（就结构而言：$p^2 - 1 = (p-1)(p+1)$ 是
两个相邻偶数之积，故 $8 \mid p^2 - 1$；
循环群 $\mathbb F_{p^2}^\times$（循环性证明见
\cref{ch:b3:galois}）则含元素 $\zeta$，其阶为 $8$，
是 $X^4 + 1$ 的根；其在 $\mathbb
F_p$ 上的极小多项式整除 $X^4+1$，且次数 $\leq 2$——$X^4+1$ 绝不
可能模 $p$ 不可约。）
\end{solution}

\begin{solution}{exo:b3:rings:12}
(a) $(1-e)^2 = 1 - 2e + e^2 = 1 - e$。映射 $\varphi(x) =
(ex, (1-e)x)$ 是加法的，且乘法地映入两个\omterm{def:b3:rings:ideal}{理想}之积：
$exey = e^2xy = e(xy)$，且 $Ae$ 是一个
交换环，单位元为 $e$（$e\cdot ex = ex$）。单射：
$ex = 0$ 与 $(1-e)x = 0$ 相加得 $x = 0$。满射：$(ea,
(1-e)b)$ 是 $ea + (1-e)b$ 的像（两个分量均可利用
$e(1-e) = 0$ 计算）。单位元正确地映为 $(1, 0)$ 型
的数对：$\varphi(1) = (e, 1-e)$，即乘积的
单位元。

(b) 在整环中，$e(e - 1) = 0$ 迫使 $e \in \{0, 1\}$：只有
平凡幂等元。在 $\Z/12\Z$ 中解 $e^2 \equiv e$：$e
\in \{0, 1, 4, 9\}$。非平凡的一对 $\{4, 9\}$：$4 + 9 =
13 \equiv 1$，$4\cdot9 = 36 \equiv 0$，且 $\Z/12\Z\cdot4 =
\{0, 4, 8\} \cong \Z/3\Z$（单位元 $4$），$\Z/12\Z\cdot9 = \{0,
3, 6, 9\} \cong \Z/4\Z$（单位元 $9$）：中国剩余定理给出的分解
$\Z/12\Z \cong \Z/4\Z\times\Z/3\Z$，其中 $9 \leftrightarrow
(1, 0)$，$4 \leftrightarrow (0, 1)$。

(c) 在中国剩余定理的同构 $\Z/n\Z \cong \prod_i
\Z/p_i^{a_i}\Z$ 下，满足所述同余条件的元素 $e_i$
对应于 $1$ 位于第 $i$ 个位置、其余位置为
$0$ 的元组，即积环的基本幂等元。
反之，一族完备的正交幂等元
（$e_ie_j = 0$，$i \neq j$；$\sum e_i = 1$）
由 (a) 归纳地还原出积的分解。
幂等元之于环，正如正交投影之于希尔伯特空间
（\cref{ch:b3:hilbert}）：内直和
分解的坐标。
\end{solution}

\begin{solution}{pb:b3:rings:1}
\textbf{1.} 有 $N(z) = z\bar z$，于是 $N(zw) = zw\overline{zw} = z\bar
z\, w \bar w = N(z)N(w)$。若 $uv = 1$，则 $N(u)N(v) = 1$（在 $\N$ 中），故
$N(u) = 1$，即 $u \in \{\pm 1, \pm \iu\}$；此四者均为单位元。
婆罗摩笈多：$(a^2+b^2)(c^2+d^2) = N\bigl((a + \iu b)(c + \iu
d)\bigr) = (ac - bd)^2 + (ad + bc)^2$.

\textbf{2.} 若 $z = ab$，则 $N(z) = N(a)N(b)$ 为质数，故
$N(a) = 1$ 或 $N(b) = 1$：其中一个因子为单位元。由于 $N(z) > 1$，
$z$ 既非零也非单位元，是\omterm{def:b3:rings:divisibility}{不可约元}——并且是\omterm{def:b3:rings:divisibility}{素元}，因为
$\Z[\iu]$ 是\omterm{def:b3:rings:pidufd}{唯一分解整环}
(\cref{thm:b3:rings:euclideanpid,thm:b3:rings:pidufd,%
lem:b3:rings:bezout}).

\textbf{3.} $\pi$ 整除 $N(\pi) = \pi\bar\pi \geq 2$，它是一个
整数；将 $N(\pi)$ 分解为质数之积，并利用
$\pi$ 为\omterm{def:b3:rings:divisibility}{素元}，故有 $\pi \mid p$，其中 $p$ 为某个质数。若又
$\pi \mid q \neq p$，则在 $\Z$ 中由贝祖定理得 $1 = up + vq$，故
$\pi \mid 1$——矛盾，故 $p$ 唯一。由 $p = \pi\gamma$：
$p^2 = N(p) = N(\pi)N(\gamma)$，且 $N(\pi) \neq 1$，故 $N(\pi)
\in \{p, p^2\}$.

\textbf{4.} 设 $\pi$ 为整除 $p$ 的高斯\omterm{def:b3:rings:divisibility}{素元}，$p =
\pi\gamma$。若 $N(\pi) = p^2$，则 $N(\gamma) = 1$，故 $p$ 为
$\pi$ 的\omterm{def:b3:rings:divisibility}{相伴}元，本身即高斯\omterm{def:b3:rings:divisibility}{素元}；且不存在
范数为 $p$ 的高斯\omterm{def:b3:rings:divisibility}{素元}（若 $N(\rho) = p$，则 $\rho \mid
\rho\bar\rho = p$，而 $p$ 在 $\Z[\iu]$ 中为\omterm{def:b3:rings:divisibility}{素元}，将迫使
$\rho$ 与 $p$ \omterm{def:b3:rings:divisibility}{相伴}，从而 $N(\rho) = N(p) = p^2 \neq
p$).
若 $N(\pi) = p$：记 $\pi = a + \iu b$，则 $p = \pi\bar\pi = a^2
+ b^2$.

\textbf{5.} 在阿贝尔群 $(\Z/p\Z)^\times$ 中，将每个
元素与其逆元配对。自逆元是
$X^2 - 1$ 的根：恰为 $\pm 1$（域中至多两个根）。
于是全体元素之积为 $1 \cdot (-1) \cdot \prod
(\text{ 配对 } k k^{-1}) = -1$：$(p-1)! \equiv -1 \pmod p$。（对于
$p = 2$ 时：$1! \equiv -1$。）

\textbf{6.} 记 $(p-1)! = \prod_{k=1}^m k \cdot
\prod_{k=m+1}^{p-1}k$，其中 $m = \frac{p-1}2$。在第二个
乘积中作代换 $k = p - j$，$j = 1, \dots, m$，在模 $p$ 下
$\prod_{j=1}^m (p - j) \equiv (-1)^m m!$。于是 $-1 \equiv
(-1)^m (m!)^2$，即 $(m!)^2 \equiv (-1)^{m+1} \pmod p$。

\textbf{7.} 若 $p \equiv 1 \pmod 4$，则 $m$ 为偶数，由第 6 题
得 $(m!)^2 \equiv -1$，即 $-1$ 的一个平方根。反之，若
$x^2 \equiv -1 \pmod p$（$p$ 为奇数），则 $x^4 = 1 \neq x^2$：$x$
的阶为 $4$（在 $(\Z/p\Z)^\times$ 中），故 $4 \mid p - 1$（由拉格朗日定理）。
又当 $p = 2$ 时：$1^2 = 1 \equiv -1$。
结论：$-1$ 为模 $p$ 的平方，当且仅当 $p = 2$ 或 $p \equiv 1 \pmod 4$。

\textbf{8.} 由 $x^2 \equiv -1$ 得 $p \mid x^2 + 1 = (x + \iu)(x -
\iu)$。若 $p$ 为高斯\omterm{def:b3:rings:divisibility}{素元}，则它须整除某一因子；但
$\frac xp \pm \frac \iu p \notin \Z[\iu]$。故 $p$ 不是
高斯\omterm{def:b3:rings:divisibility}{素元}；由二分性（第 4 题）——$p$ 非\omterm{def:b3:rings:divisibility}{素元}即意味着
第二分支，故 $p = a^2 + b^2$。

\textbf{9.} 平方数 $\equiv 0$ 或 $1 \pmod 4$，故 $a^2 + b^2
\in \{0, 1, 2\} \pmod 4$：素数 $p \equiv 3 \pmod 4$ 不可
表为两个平方数之和。由第 4 题，分支 $N(\pi) = p$（$p =
a^2+b^2$）不可能：$p$ 仍为高斯\omterm{def:b3:rings:divisibility}{素元}。

\textbf{10.} $(1 + \iu)^2 = 2\iu$，故 $2 = -\iu(1 + \iu)^2$；且
$N(1 + \iu) = 2$ 为素数，故 $1 + \iu$ 为高斯\omterm{def:b3:rings:divisibility}{素元}
（第 2 题）。

\textbf{11.} 每个高斯\omterm{def:b3:rings:divisibility}{素元}恰整除一个素数
$p$（第 3 题）；按情形列出：$p = 2$ 时给出
$1 + \iu$ 的\omterm{def:b3:rings:divisibility}{相伴}元；$p \equiv 3 \pmod 4$ 时给出 $p$ 本身
（第 9 题）；$p \equiv 1 \pmod 4$ 时给出一对 $\pi, \bar\pi$
范数为 $p$（第 4、8 题）。这一对确实不同：$\bar\pi \in
\{\pm\pi, \pm\iu\pi\}$ 将迫使，记 $\pi = a + \iu b$，
或 $b = 0$、$a = 0$ 或 $a = \pm b$，由此 $p = a^2 + b^2 \in
\{a^2, 2a^2\}$ —— 对奇素数不可能。验算：
$5 = (2 + \iu)(2 - \iu)$，$N(2\pm\iu) = 5$；$3$ 为范数
$9$ 的\omterm{def:b3:rings:divisibility}{素元}。

\textbf{12.} 记 $n = 2^{\alpha}\prod_i p_i^{\beta_i} \prod_j
q_j^{2\gamma_j}$，其中 $p_i \equiv 1$，$q_j \equiv 3 \pmod 4$。
每个因子都是两个平方数之和：$2 = 1^2 + 1^2$；$p_i = a^2
+ b^2$（第 8 题）；$q_j^{2\gamma_j} = (q_j^{\gamma_j})^2 +
0^2$。婆罗摩笈多恒等式（第 1 题）将
该性质传递到乘积 $n$ 上。

\textbf{13.} 设 $n = a^2 + b^2 = N(a + \iu b)$，$q \equiv 3
\pmod 4$，$q \mid n$。高斯\omterm{def:b3:rings:divisibility}{素元} $q$（第 9 题）整除
$(a + \iu b)(a - \iu b)$，故整除两个因子之一——不妨设 $q
\mid a + \iu b$（另一情形完全类似）。但此时 $\frac{a +
\iu b}{q} = \frac aq + \iu \frac bq \in \Z[\iu]$ 表明 $q
\mid a$ \emph{且} $q \mid b$（在 $\Z$ 中）。故 $q^2 \mid n$，并且
$\frac n{q^2} =
\bigl(\frac aq\bigr)^2 + \bigl(\frac bq\bigr)^2$。由
对 $n$ 的强归纳法，$q$ 在 $n/q^2$ 中的指数为偶数；
$n$ 中该指数同样为偶数。

\textbf{14.} \emph{定理（费马）。}正整数可表为两个平方数之和，当且仅当
每个 $\equiv 3 \pmod 4$ 的素数
在其中以偶数次幂出现。——$2025 = 3^4 \cdot 5^2$：
$3$ 的指数为偶数，是（$2025 = 45^2 + 0^2 = 27^2 + 36^2$）。
$2026 = 2 \cdot 1013$，$1013 \equiv 1 \pmod 4$ 为素数：是
（$1013 = 22^2 + 23^2$，由婆罗摩笈多恒等式及 $2 = 1^2+1^2$：
$2026 = (22 - 23)^2 + (22 + 23)^2 = 1^2 + 45^2$）。$2027$ 为
$\equiv 3 \pmod 4$ 的素数：否。

\textbf{15.} 设 $p = a^2 + b^2 = c^2 + d^2$，其中正整数
$p \equiv 1 \pmod 4$，且 $\pi$ 为高斯\omterm{def:b3:rings:divisibility}{素元}，满足
$p = \pi\bar\pi$（第 4 题）。$a + \iu b$ 与 $c + \iu d$ 的范数均为 $p$，
因而均为高斯\omterm{def:b3:rings:divisibility}{素元}（第 2 题），且整除
$p = (a+\iu b)(a - \iu b)$。由分解的唯一性，$c +
\iu d$ 为 $a + \iu b$ 或 $a - \iu b$ 的\omterm{def:b3:rings:divisibility}{相伴}元：
\[
c + \iu d \in \{\pm(a \pm \iu b),\ \pm\iu(a \pm \iu b)\}
= \{\pm a \pm \iu b,\ \pm b \pm \iu a\}.
\]
由 $c, d$ 为正，只余 $c + \iu d \in \{a + \iu b, b + \iu
a\}$，即 $\{c, d\} = \{a, b\}$。

\textbf{16.} 若 $n = a^2 - b^2 = (a-b)(a+b)$：两个因子
奇偶性相同，故 $n$ 为奇数（二者皆奇）或被
$4$ 整除（二者皆偶）——绝不 $\equiv 2 \pmod 4$。反之，$n$
为奇数：$n = \bigl(\frac{n+1}2\bigr)^2 -
\bigl(\frac{n-1}2\bigr)^2$；$n = 4m$：$n = (m+1)^2 - (m-1)^2$。
答案是一个裸的同余条件，背后只需一行
恒等式：平方差不含算术上的
深度，而和的对比正是这道题的全部意义
所在。

\textbf{17.} $(a, b) \mapsto z = a + \iu b$ 是表示与
$\{z : N(z) = n\}$ 之间的双射。分解 $z$ 时，在
\omterm{def:b3:rings:pidufd}{唯一分解整环} $\Z[\iu]$ 中，利用分类（第 11 题），相差
一个单位，$z = (1+\iu)^{a'}\prod_j\pi_j^{s_j}\bar\pi_j^{t_j}
\prod_kq_k^{u_k}$，并取范数（$N(1+\iu) = 2$，
$N(\pi_j) = N(\bar\pi_j) = p_j$，$N(q_k) = q_k^2$）：
\[
n = 2^{a'}\prod_jp_j^{s_j + t_j}\prod_kq_k^{2u_k} .
\]
比较指数：$a' = a$，$s_j + t_j = b_j$，$2u_k = c_k$
——当且仅当每个 $c_k$ 为偶数时有解，此时 $u_k = c_k/2$
被唯一确定，而 $s_j \in \intint0{b_j}$ 可自由选取。不同的数据
$(u, (s_j))$ 给出互不\omterm{def:b3:rings:divisibility}{相伴}的 $z$，它们具有\emph{相同}的
范数；单位 $u \in \{\pm1, \pm\iu\}$（共 4 种选择）则
逐一枚举每个\omterm{def:b3:rings:divisibility}{相伴}类，且不重复（两个相等的
乘积将违背分解的唯一性——
$\pi_j$ 与 $\bar\pi_j$ 互不\omterm{def:b3:rings:divisibility}{相伴}，因为 $p_j =
\pi_j\bar\pi_j$ 不分歧）。总计 $r_2(n) =
4\prod_j(b_j + 1)$，而 $0$ 出现于某个 $c_k$ 为奇数的情形。

\textbf{18.} $\chi(dd') = \chi(d)\chi(d')$ 在模 $4$ 下验证
（奇数 $\times$ 奇数涵盖四种符号情形；凡含偶数因子者
皆得 $0 = 0$）。设 $m, n$ 互素，则 $mn$ 的因数
可唯一地写作 $d = d_1d_2$，其中 $d_1 \mid m$，$d_2 \mid n$：
$\sum_{d \mid mn}\chi(d) = \bigl(\sum_{d_1\mid
m}\chi(d_1)\bigr)\bigl(\sum_{d_2\mid n}\chi(d_2)\bigr)$:
故为乘性。素数幂：在 $2^a$ 上，仅 $d = 1$ 为奇数：
和 $= 1$。在 $p^b$（$p \equiv 1$）上，所有 $\chi(p^i) = 1$，
和 $= b + 1$。在 $q^c$（$q \equiv 3$）上：$\chi(q^i) =
(-1)^i$，交错和 $= 1$（$c$ 为偶数）或 $0$（$c$ 为奇数）。

\textbf{19.} 两个乘性函数
$\frac14r_2$（第 17 题）与 $\sum_{d\mid n}\chi(d)$
（第 18 题）在一切素数幂上相等——取值 $1$（$2^a$ 处）；
$b + 1$（$p^b$ 处）；$\mathbf 1_{c\ \mathrm{even}}$（$q^c$ 处）——
因此处处相等，即雅可比公式，其中 $\sum_{d\mid
n}\chi(d) = d_1(n) - d_3(n)$ 由因数分类得出。验算：
$r_2(3) = 0 = 4(1 - 1)$; $r_2(5) = 8 = 4(2 - 0)$
($(\pm1,\pm2), (\pm2,\pm1)$); $r_2(9) = 4 = 4(2 - 1)$
（因数 $1, 9 \equiv 1$；$3 \equiv 3$；表示
$(\pm3, 0), (0, \pm3)$）；$r_2(25) = 12 = 4(3 - 0)$。对 $65 =
5\cdot13$：$r_2 = 4\cdot2\cdot2 = 16$，由 $65 = 1 + 64 = 16
+ 49$：十六组序对为 $(\pm1, \pm8), (\pm8, \pm1), (\pm4,
\pm7), (\pm7, \pm4)$.

\textbf{20.} $\sum_{n \leq x}r_2(n)$ 计数序对 $(a, b)$
满足 $0 < a^2 + b^2 \leq x$ 者，即
闭圆盘 $D_{\sqrt x}$ 去掉原点后的格点。对每个
格点 $P$ 配以单位正方形 $P + \intco01^2$：这些
正方形铺满平面。凡附于
$D_{\sqrt x}$ 中某点的正方形皆落在 $D_{\sqrt x + \sqrt2}$ 内，而
与 $D_{\sqrt x - \sqrt 2}$ 相交的正方形皆附于
$D_{\sqrt x}$ 中某点（正方形直径为 $\sqrt 2$）：比较
面积，有
\[
\pi(\sqrt x - \sqrt2)^2 \leq \#\{\text{格点，位于 }
D_{\sqrt x}\} \leq \pi(\sqrt x + \sqrt 2)^2,
\]
两端均为 $\pi x + O(\sqrt x)$。减去
原点在此精度下无影响。

\textbf{21.} 由雅可比公式（第 19 题）并交换
求和次序（$n = dm$）：
\[
\frac14\sum_{n\leq x}r_2(n) = \sum_{n \leq x}\sum_{d \mid
n}\chi(d) = \sum_{d \leq x}\chi(d)\,\#\{m : dm \leq x\}
= \sum_{d\leq x}\chi(d)\Bigl\lfloor\frac xd\Bigr\rfloor,
\]
其值为 $\frac{\pi x}4 + O(\sqrt x)$（由第 20 题）。去掉
取整：$\lfloor x/d\rfloor = x/d + O(1)$，但
对 $O(1)$ 在 $d \leq x$ 上求和过于粗糙；改用
$\chi$ 的部分和有界（$0, 1, 1, 0$ 循环），
故由阿贝尔求和，$\sum_{d\leq x}\chi(d)\{x/d\}$ 的
项两两分组（$d \equiv 1, 3$），其和为 $O(\sqrt x)$——
另法更简单：在 $\sqrt x$ 处分段。对 $d \leq
\sqrt x$，将 $\lfloor x/d\rfloor$ 换作 $x/d + O(1)$：误差
为 $O(\sqrt x)$。对 $d > \sqrt x$，$\lfloor x/d\rfloor$ 的
每个取值 $v < \sqrt x$ 都在一段连续的 $d$ 上，
其上 $\chi$ 之和为 $O(1)$：总误差为 $O(\sqrt x)$，由
对 $\leq \sqrt x$ 个 $v$ 的取值求和得之；而
$\sum_{d > \sqrt x}\chi(d)\frac xd = O(\sqrt x)$，由
交错级数尾项估计（$x\sum_{d>\sqrt x}\chi(d)/d =
x\,O(1/\sqrt x)$）。于是
\[
x\sum_{d \leq x}\frac{\chi(d)}d = \frac{\pi x}4 + O(\sqrt x),
\qquad\text{即}\qquad
\sum_{d\leq x}\frac{\chi(d)}d = \frac\pi4 +
O\Bigl(\frac1{\sqrt x}\Bigr),
\]
令 $x \to \infty$，得 $1 - \frac13 + \frac15 - \dots =
\frac\pi4$.

\textbf{22.} 若 $n \equiv 3 \pmod4$ 可写成 $a^2 + b^2$：平方数
$\equiv 0, 1 \pmod 4$，故 $a^2 + b^2 \in \{0, 1, 2\}$
模 $4$ 下不可能——（第 17 题的判别准则所述相同：
$n \equiv 3 \pmod 4$ 迫使某个 $\equiv 3$ 的素数指数为
奇数）。因此可表示的整数避开一整个
剩余类，即其密度 $\leq \frac34$。$\pi$ 为 $r_2$ 的平均值，
它集中于少数整数上：$n = \prod_{j\leq k}p_j$
（互异素数，$\equiv 1 \bmod 4$）满足 $r_2(n) = 4\cdot2^k$
种表示——无界多——故一个稀疏的
$n$ 之集合即可承载整个平均值，恰如彩票的平均收益
与几乎必然的亏损并存。兰道的
估计 $\#\{n \leq x \text{ 可表示}\} \sim
Cx/\sqrt{\log x}$ 证实了这一点：密度为 $0$，平均值为 $\pi$。

\textbf{23.} \emph{必要性。}设 $n = a^2 + b^2$，其中
$\gcd(a, b) = 1$。若素数 $q \equiv 3 \pmod 4$ 整除 $n$，
由第 13 题可得 $q \mid a$ 且 $q \mid b$，矛盾。若
$4 \mid n$：平方数 $\equiv 0, 1 \pmod 4$，故 $a^2 + b^2
\equiv 0 \pmod 4$ 迫使 $a^2 \equiv b^2 \equiv 0$，即 $a,
b$ 皆为偶数，矛盾。\emph{充分性。}记 $n =
2^{\alpha}\prod_jp_j^{b_j}$，其中 $\alpha \leq 1$ 且 $p_j
\equiv 1 \pmod 4$，并令 $z = (1+\iu)^{\alpha}\prod_j
\pi_j^{b_j} = a + \iu b$，其范数为 $n$。设素数 $t$
整除 $\gcd(a, b)$，则 $t \mid z$（在 $\Z[\iu]$ 中）。若 $t
\equiv 3 \pmod 4$：$t \mid N(z) = n$，已排除。若 $t \equiv 1
\pmod 4$：$t = \pi_t\bar\pi_t$，故 $\bar\pi_t \mid z$；但
$z$ 的分解式中不含共轭\omterm{def:b3:rings:divisibility}{素元}（$\pi_j$ 与
$\bar\pi_j$ 互不\omterm{def:b3:rings:divisibility}{相伴}，第 17 题），与
唯一分解相矛盾。若 $t = 2 = -\iu(1+\iu)^2$，则
$(1+\iu)^2 \mid z$，迫使 $\alpha \geq 2$，已排除。故
$\gcd(a, b) = 1$，即该表示是本原的。

\textbf{24.} $a$ 为奇数（$\gcd(a, b) = 1$，$b$ 为偶数），故 $c^2
= a^2 + b^2$ 为奇数，$c$ 亦为奇数。设 $\delta$ 为
$a + \iu b$ 与 $a - \iu b$ 的公共高斯素因子，它
整除其和 $2a$ 及其差 $2\iu b$，从而
$2a$ 与 $2b$ 亦被其整除，由贝祖关系 $ua + vb = 1$ 即得
$\delta \mid 2$，故 $\delta$ 与 $1 + \iu$ \omterm{def:b3:rings:divisibility}{相伴}，且
$N(\delta) = 2$ 整除 $N(a + \iu b) = c^2$，而后者为奇数：
矛盾。故 $a + \iu b$ 与 $a - \iu b$ 互素，
乘积为 $c^2$；在\omterm{def:b3:rings:pidufd}{唯一分解整环} $\Z[\iu]$ 中，
$c^2$ 的每个高斯\omterm{def:b3:rings:divisibility}{素元}以偶数次幂出现，并整个落入
两个互素因子之一，故 $a + \iu b = u(m + \iu
n)^2 = u\bigl(m^2 - n^2 + 2\iu mn\bigr)$，其中 $u$ 为单位。
$u = \pm\iu$ 的情形使实部 $\mp 2mn$ 为偶数——
不可能，因 $a$ 为奇数。$u = \pm1$ 的取法给出，
调整 $m, n$ 的符号并交换其名字，使
一切为正，则 $a = m^2 - n^2$，$b = 2mn$，其中 $m > n
\geq 1$；由 $c^2 = N(m + \iu n)^2$ 得 $c = m^2 + n^2$。
$m$ 与 $n$ 的公因数必整除 $a$ 与 $b$：$\gcd
(m, n) = 1$；若 $m \equiv n \pmod 2$，则 $a$ 为偶数：
奇偶性相反。验算：$(m, n) = (2, 1)$ 给出 $(3, 4,
5)$；$(m, n) = (5, 2)$ 给出 $(25 - 4, 20, 25 + 4) = (21, 20,
29)$，且 $441 + 400 = 841 = 29^2$。

\textbf{25.} 由雅可比公式 $r_2(n) = 4(d_1(n) - d_3(n))$
可得，对 $n = 1, \dots, 25$：
\[
4,\ 4,\ 0,\ 4,\ 8,\ 0,\ 0,\ 4,\ 4,\ 8,\ 0,\ 0,\ 8,\ 0,\ 0,\
4,\ 8,\ 4,\ 0,\ 8,\ 0,\ 0,\ 0,\ 0,\ 12,
\]
非零恰在 $n = 1, 2, 4, 5, 8, 9, 10, 13, 16, 17, 18,
20, 25$ 处（例如 $r_2(15) = 0$：因数 $1, 5 \equiv 1$
与 $3, 15 \equiv 3$ 相抵；$r_2(20) = 8$：因数 $1, 5
\equiv 1$，无 $\equiv 3$ 者）。总和为 $4 + 4 + 4 + 8 + 4 +
4 + 8 + 8 + 4 + 8 + 4 + 8 + 12 = 80$。因数一侧：
奇数 $d \leq 25$ 的贡献为
\[
25 - 8 + 5 - 3 + 2 - 2 + 1 - 1 + 1 - 1 + 1 - 1 + 1 = 20,
\]
即取 $\chi(d)\lfloor 25/d\rfloor$，$d = 1, 3, 5, \dots,
25$；$4 \cdot 20 = 80$，与第 21 题的
恒等式所预言者一致。半径为 $5$ 的闭圆盘上的格点：
$80$ 个满足 $1 \leq a^2 + b^2 \leq 25$ 的点，外加
原点，即 $81$；而 $\pi x = 25\pi \approx 78.54$，误差
约为 $2.46$，完全落在第 20 题所给的 $O(\sqrt x)$
范围之内（$\sqrt x = 5$）。
\end{solution}
